\documentclass[leqno]{article}
% Shared QBT/Soln pack preamble for resources.danieljsmith.org.
\usepackage{graphicx}
\usepackage{dsfont}
\usepackage{mathtools}
\usepackage{amssymb}
\usepackage{amsmath}
\usepackage{mleftright}
\usepackage{amsthm}
\usepackage{float}
\usepackage{ragged2e}
\usepackage{tensor}
\usepackage{fancyhdr}
\usepackage{empheq}
\usepackage[most]{tcolorbox}
\usepackage{enumitem}
\usepackage{dirtytalk}
\usepackage{marginnote}
\usepackage{microtype}
\usepackage[
  a4paper,
  left=0.8in,
  right=1in,
  top=1in,
  bottom=0.5in,
  headheight=14pt,
  headsep=0.4in
]{geometry}
\setlength{\parindent}{0cm}

% TikZ / PGFPlots
\usepackage{pgfplots}
\pgfplotsset{compat=1.18}
\usepgfplotslibrary{fillbetween, polar}
\usetikzlibrary{calc, positioning}
\usetikzlibrary{angles, quotes}
\usetikzlibrary{arrows.meta}
\usetikzlibrary{decorations.pathreplacing, decorations.markings}
\usetikzlibrary{patterns}
\usetikzlibrary{intersections}

% Links and contents. The contents follow the TMUA papers: compact entries,
% no dot leaders, and each question's mark total in grey after its label.
\usepackage{tocloft}
\tocloftpagestyle{fancy}
\setlength{\cftbeforesubsecskip}{2pt}
\setlength{\cftsubsecindent}{1.5em}
\setlength{\cftsubsecnumwidth}{0pt}
\renewcommand{\cftsubsecfont}{\small\color{black}}
\renewcommand{\cftsubsecpagefont}{\small\color{black}}
\renewcommand{\cftsubsecleader}{\hfill}
\usepackage[colorlinks=true,linkcolor=black,citecolor=magenta,urlcolor=black]{hyperref}%
\urlstyle{same}

% Pack macros
% \marks also adds its value to the current question's total in the .aux
% file, so the contents show each total from the second pass. The write
% happens at shipout, so a box typeset more than once counts once. The
% pack's mark total for the contents-page title is summed the same way.
\makeatletter
\newcount\djs@totalmarks
\newcommand{\djsAddMarks}[2]{%
  \global\advance\djs@totalmarks#2\relax
  \@ifundefined{djs@marks@#1}%
    {\expandafter\xdef\csname djs@marks@#1\endcsname{#2}}%
    {\expandafter\xdef\csname djs@marks@#1\endcsname
      {\the\numexpr\csname djs@marks@#1\endcsname+#2\relax}}}
\RenewDocumentCommand{\marks}{o m}{%
  \protected@write\@auxout{}{\string\djsAddMarks{\number\value{enumi}}{#2}}%
  \IfValueTF{#1}%
    {\marginnote{\raggedright\textbf{[#2]}}[#1]}%
    {\marginnote{\raggedright\textbf{[#2]}}}%
}
\makeatother

\newcommand{\vspacerule}[2]{%
  \par\vspace*{#1}%
  \noindent\hrulefill\par
  \vspace*{#2}%
}

\definecolor{scarlet}{RGB}{205, 40, 75}
% Site colours for the contents-page title: the accent, the rule colour, a
% stronger rule colour for the path separator, and the muted text colour.
\definecolor{djsBurgundy}{RGB}{118, 59, 67}
\definecolor{djsHairline}{RGB}{216, 213, 205}
\definecolor{djsSlash}{RGB}{170, 166, 158}
\definecolor{djsMuted}{RGB}{96, 101, 106}

% Worked-solution box, adapted from the TMUA worked-solutions box. A white breakable box with a left scarlet rule and no surrounding
% frame or tint. A part that breaks fills the rest of its page, so the rule
% runs to the bottom margin; the unbroken or last part ends in a short
% horizontal foot rule so a solution visibly stops. Unlike the TMUA box there
% is no answer label. The box does not grow to the left, so the rule sits at
% the text edge; growing to the right keeps the usable line width close to the
% old tcolorbox Solution box.
% Spacing inside the box: paragraphs are 6pt apart, so source blank lines show
% (solutions put one after each display and none between a lead-in and its
% display). Every display gets 6pt above and 10pt below, whatever the length
% of its lead-in line; align rows are spread by an 8pt \jot; solution parts
% are 14pt apart and sub-parts 8pt; \left and \right space like ordinary
% brackets (mleftright); no single line is stranded at a page break; and
% neither the Solution label nor the line before a list ends a page. No
% skip can shrink, so full pages are not squeezed.
\newtcolorbox{djsSolution}{%
  enhanced, breakable, frame hidden, sharp corners, colback=white,
  height fixed for=first and middle, height fill,
  extras unbroken and last={natural height},
  borderline west={2.2pt}{0pt}{scarlet},
  left=11pt, right=0pt, top=5pt, bottom=13pt,
  before skip=13pt, after skip=4pt,
  grow to right by=10mm,
  overlay unbroken and last={%
    \fill[scarlet] (frame.south west)
      rectangle ([xshift=3.5cm, yshift=2.2pt]frame.south west);},
  before upper={{\color{scarlet}\bfseries Solution}\par\nobreak\vspace{4pt}%
    \setlength{\parskip}{6pt plus 2pt}%
    \setlength{\abovedisplayskip}{6pt plus 2pt}%
    \setlength{\abovedisplayshortskip}{6pt plus 2pt}%
    \setlength{\belowdisplayskip}{10pt plus 2pt}%
    \setlength{\belowdisplayshortskip}{10pt plus 2pt}%
    \setlength{\jot}{8pt}%
    \setlist{parsep=6pt plus 2pt, topsep=6pt plus 2pt, beginpenalty=10000}%
    \setlist[enumerate,1]{itemsep=14pt plus 2pt}%
    \setlist[enumerate,2]{itemsep=8pt plus 2pt}%
    \mleftright
    \clubpenalty=10000 \widowpenalty=10000 \displaywidowpenalty=10000\relax}}

% Optional smaller figure inside a solution box. \scalebox shrinks node text
% with the drawing. Not applied to existing figures.
\newcommand{\SolnFigureScale}{0.85}
\newsavebox{\djsSolnFigureBox}
\newenvironment{djsSolutionFigure}
  {\par\vspace{4pt}\begin{center}\begin{lrbox}{\djsSolnFigureBox}}
  {\end{lrbox}\scalebox{\SolnFigureScale}{\usebox{\djsSolnFigureBox}}%
   \end{center}\vspace{2pt}\par}

\newcommand{\cxl}[2]{%
  \tikz[baseline=(X.base)]{%
    \node[inner sep=0pt, outer sep=0pt, text=#1] (X) {$#2$};
    \draw[#1, line width=0.4pt] (X.south west)--(X.north east);
  }%
}

% Page-style hook run at the contents page and each question's opening page.
% A no-op for QBT packs; \djsSolnHeader points it at the fancy header.
\newcommand{\djsQuestionPageStyle}{}

% A contents entry: "Question N" in a box wide enough for two digits, so the
% mark totals line up, then the total in grey. Bookmarks keep "Question N".
\makeatletter
\newlength{\djsTocLabelWidth}
\newcommand{\djsTocMarks}[1]{%
  \@ifundefined{djs@marks@#1}{}{%
    {\footnotesize\itshape\color{black!45}%
     \csname djs@marks@#1\endcsname~mark%
     \ifnum\csname djs@marks@#1\endcsname=1 \else s\fi}}}
\newcommand{\djsTocEntry}[1]{%
  \texorpdfstring
    {\settowidth{\djsTocLabelWidth}{Question 00}%
     \makebox[\dimexpr\djsTocLabelWidth+0.9em\relax][l]{Question #1}%
     \djsTocMarks{#1}}%
    {Question #1}}
\makeatother

% The question count for the contents-page title is written to the .aux at
% the end of the document, so it also appears from the second pass.
\newcounter{djsQuestions}
\newcommand{\questionitem}{%
  \djsQuestionPageStyle
  \item\phantomsection
  \stepcounter{djsQuestions}%
  \addcontentsline{toc}{subsection}{\protect\djsTocEntry{\number\value{enumi}}}%
}
\makeatletter
\newcommand{\djsSetQuestionCount}[1]{\gdef\djs@questioncount{#1}}
\AtEndDocument{\immediate\write\@auxout
  {\string\djsSetQuestionCount{\number\value{djsQuestions}}}}
\makeatother

\renewcommand{\contentsname}{Questions}

% \djsPackHeader[<link content tag>]{<left header>}; the tag says which
% document a click on the header link came from.
\newcommand{\djsPackHeader}[2][questions]{%
  \pagestyle{fancy}%
  \fancyhead{}%
  \fancyfoot{}%
  \renewcommand{\headrulewidth}{0.9pt}%
  \lhead{#2}%
  \chead{}%
  \rhead{\href{https://danieljsmith.org/?utm_source=pdf&utm_campaign=further-maths&utm_content=#1}{danieljsmith.org}}%
}

% The header's topic is also the contents-page title, under which a label
% says which booklet this is.
\newcommand{\djsPackTitle}{}
\newcommand{\djsBookletLabel}{}
\newcommand{\djsQbtHeader}[1]{%
  \djsPackHeader{#1}%
  \renewcommand{\djsPackTitle}{#1}%
  \renewcommand{\djsBookletLabel}{Questions by topic}%
}

% The colophon at the foot of page one, independent of the page style: the
% contents-page title's hairline led by a short scarlet rule, then the date on
% the left and a link to the site on the right, all inside the bottom margin.
% Its sizes are fixed, so it matches the TMUA papers' colophon. The document
% supplies the date; recompiling never changes it.
\makeatletter
\newcommand{\djsLastUpdated}[1]{%
  \AtBeginDocument{\AddToHookNext{shipout/foreground}{\djs@colophon{#1}}}}
\newcommand{\djs@colophon}[1]{%
  \put(0.8in,-\dimexpr\paperheight-0.47in\relax){%
    \makebox[0pt][l]{\color{djsHairline}%
      \rule[0.9pt]{\dimexpr\paperwidth-1.8in\relax}{0.6pt}}%
    \makebox[0pt][l]{\color{scarlet}\rule{3.5cm}{2.2pt}}}%
  \put(0.8in,-\dimexpr\paperheight-0.32in\relax){\makebox[0pt][l]{%
    \fontsize{9}{11}\selectfont Last updated #1}}%
  \put(\dimexpr\paperwidth-1in\relax,-\dimexpr\paperheight-0.32in\relax){%
    \makebox[0pt][r]{\fontsize{7}{8.5}\selectfont\color{djsMuted}%
      Latest version at \href{https://resources.danieljsmith.org/}%
        {\color{black}resources.danieljsmith.org}}}}
\makeatother

% Solutions packs show the header only on the contents page and each
% question's opening page, so it never appears part-way through a solution.
% \thispagestyle marks the next page shipped out, which is the question's
% page because every solution question starts after a \newpage.
\newcommand{\djsSolnHeader}[1]{%
  \djsPackHeader[solutions]{\textit{Solutions} - #1}%
  \pagestyle{empty}%
  \renewcommand{\djsQuestionPageStyle}{\thispagestyle{fancy}}%
  \renewcommand{\djsPackTitle}{#1}%
  \renewcommand{\djsBookletLabel}{Worked solutions}%
}

% \djsSitePath{<strand>}{<topic>}, called once after the header, names where
% the pack sits on the site; the contents-page title prints it above the pack
% title, e.g. CORE PURE / FURTHER CALCULUS. Without it the line is left out.
\makeatletter
\newif\ifdjs@sitepath
\newcommand{\djs@strand}{}
\newcommand{\djs@section}{}
\newcommand{\djsSitePath}[2]{%
  \djs@sitepathtrue
  \renewcommand{\djs@strand}{#1}%
  \renewcommand{\djs@section}{#2}%
}
\makeatother

% Further Mechanics packs call \djsMechanicsPack in their preamble, so that
% \djsFrontMatter prints the default value of g above the contents. The box
% keeps its own space above it (before skip is the whole gap below the title,
% not an addition to it). Without the box, \djsGravityNote is just extra space
% before the contents heading, bringing the gap above it close to the gap
% between the box and the heading on Further Mechanics pages.
\newcommand{\djsGravityNote}{\vspace{14pt}}
\newcommand{\djsMechanicsPack}{%
  \renewcommand{\djsGravityNote}{%
    \begin{tcolorbox}[sharp corners, colback=white, colframe=black,
      boxrule=1pt, left=8pt, right=8pt, top=6pt, bottom=6pt,
      before skip=20pt]
    {\bfseries Instructions}\par\vspace{4pt}
    Unless a question states otherwise, take $g = 9.8\,\mathrm{m\,s^{-2}}$
    wherever a numerical value of $g$ is required. Solutions also give final
    answers for $g = 10\,\mathrm{m\,s^{-2}}$ where these differ.
    \end{tcolorbox}%
    \vspace{16pt}}%
}

% Contents-page title: the site path in spaced burgundy capitals, the pack
% title, a hairline led by a short scarlet rule (the solution box's foot rule),
% then the booklet label with the question and mark counts. The counts come
% from the .aux, so they appear from the second pass.
\makeatletter
\newsavebox{\djs@eyebrowbox}
\newcommand{\djs@caps}[1]{\textls[180]{\MakeUppercase{#1}}}
\newcommand{\djs@eyebrow}{%
  \color{djsBurgundy}\footnotesize\bfseries
  \djs@caps{\djs@strand}%
  {\color{djsSlash}\hspace{0.55em}/\hspace{0.55em}}%
  \djs@caps{\djs@section}}
\newcommand{\djs@dot}{\hspace{0.45em}\textperiodcentered\hspace{0.45em}}
\newcommand{\djs@counts}{%
  \@ifundefined{djs@questioncount}{}{%
    \djs@questioncount~question\ifnum\djs@questioncount=1 \else s\fi
    \ifnum\djs@totalmarks>0
      \djs@dot\the\djs@totalmarks~mark\ifnum\djs@totalmarks=1 \else s\fi
    \fi}}
\newcommand{\djsContentsTitle}{%
  \ifdjs@sitepath
    \sbox\djs@eyebrowbox{\djs@eyebrow}%
    \ifdim\wd\djs@eyebrowbox>\linewidth
      \@latex@warning@no@line{The \string\djsSitePath\space line is wider
        than the text, so it wraps. Shorten its topic}%
    \fi
    {\raggedright\djs@eyebrow\par}%
    \vspace{7pt}%
  \fi
  {\raggedright\fontsize{24.88}{28}\selectfont\djsPackTitle\par}%
  \vspace{9pt}%
  \noindent\makebox[0pt][l]{\color{djsHairline}\rule[0.9pt]{\linewidth}{0.6pt}}%
  {\color{scarlet}\rule{3.5cm}{2.2pt}}\par
  \vspace{5pt}%
  {\small\color{djsMuted}\djsBookletLabel\hfill\djs@counts\par}%
  \vspace{8pt}%
}
\makeatother

\newcommand{\djsFrontMatter}{%
  \djsQuestionPageStyle
  \vspace*{6pt}%
  \djsContentsTitle
  \djsGravityNote
  \tableofcontents
  \newpage
}


\djsSolnHeader{Improper Integrals and Mean Values}
\djsSitePath{Core Pure}{Further Calculus}
\djsLastUpdated{27 October 2026}

\begin{document}
\djsFrontMatter
\begin{enumerate}[
  label=\textbf{\arabic*.},
  leftmargin=*,
  topsep=0pt,
  partopsep=0pt,
  parsep=0pt
]

% ---- Question 1 ----
\questionitem
% [Source: _QBT___Solns__Improper_Integrals_and_Mean_Values, Q1]

The function $\mathrm{f}$ is defined by
\[\mathrm{f}(x) = \dfrac{x}{x^2+1}, \qquad 1 \le x \le 3\] \\
Find the mean value of $\mathrm{f}$ over the interval $1 \le x \le 3$. \\

Give your answer in the form $p\ln q$, where $p$ is a rational number and $q$ is an integer. \marks{4} \\

\vspace*{10pt}

\begin{djsSolution}
The mean value of \(\mathrm{f}(x)\) over \(a\le x\le b\) is
\[
\dfrac{1}{b-a}\int_a^b \mathrm{f}(x)\,\mathrm{d}x
\]

So here the mean value is
\[
\frac{1}{3-1}\int_1^3 \frac{x}{x^2+1}\,\mathrm{d}x
\]

The numerator is a multiple of the derivative of the denominator, since \(\frac{\mathrm{d}}{\mathrm{d}x}\left(x^2+1\right)=2x\). Also \(x^2+1>0\) for all \(x\), so the logarithm needs no modulus signs. Hence
\begin{align*}
\int_1^3 \frac{x}{x^2+1}\,\mathrm{d}x &= \frac12\int_1^3 \frac{2x}{x^2+1}\,\mathrm{d}x\\
&= \left[\dfrac12\ln\left(x^2+1\right)\right]_1^3\\
&= \dfrac12\ln 10-\dfrac12\ln 2\\
&= \dfrac12\ln 5
\end{align*}

Therefore the mean value is
\[
\frac12\times\frac12\ln 5=\frac14\ln 5
\]

This is in the form \(p\ln q\) with \(p=\frac14\) and \(q=5\).
\end{djsSolution}

\newpage

% ---- Question 2 ----
\questionitem
% [Source: _QBT___Solns__Improper_Integrals_and_Mean_Values, Q2]

Show that
\[\int_{2}^{\infty} \dfrac{x}{(x^{2}+5)^{\frac{3}{2}}}\,\mathrm{d}x = \dfrac{1}{3}\]
You must show the limiting process clearly. \marks{5} \\

\vspace*{10pt}

\begin{djsSolution}
Replace the infinite upper limit by \(t\) and take the limit:
\[
\int_{2}^{\infty} \dfrac{x}{(x^{2}+5)^{\frac{3}{2}}}\,\mathrm{d}x=\lim_{t\to\infty}\int_{2}^{t} x\left(x^{2}+5\right)^{-\frac{3}{2}}\,\mathrm{d}x
\]

The integrand is a multiple of \(f'(x)\left(f(x)\right)^{-\frac{3}{2}}\) with \(f(x)=x^2+5\), so use the reverse chain rule. Differentiating gives
\[
\frac{\mathrm{d}}{\mathrm{d}x}\left(x^{2}+5\right)^{-\frac{1}{2}}=-\dfrac{1}{2}\times 2x\left(x^{2}+5\right)^{-\frac{3}{2}}=-x\left(x^{2}+5\right)^{-\frac{3}{2}}
\]

so
\[
\int x\left(x^{2}+5\right)^{-\frac{3}{2}}\,\mathrm{d}x=-\left(x^{2}+5\right)^{-\frac{1}{2}}+c=-\dfrac{1}{\sqrt{x^{2}+5}}+c
\]

For \(t>2\):
\begin{align*}
\int_{2}^{t} \dfrac{x}{(x^{2}+5)^{\frac{3}{2}}}\,\mathrm{d}x&=\left[-\dfrac{1}{\sqrt{x^{2}+5}}\right]_{2}^{t}\\
&=-\dfrac{1}{\sqrt{t^{2}+5}}+\dfrac{1}{\sqrt{4+5}}\\
&=\dfrac{1}{3}-\dfrac{1}{\sqrt{t^{2}+5}}
\end{align*}

As \(t\to\infty\), \(\sqrt{t^{2}+5}\to\infty\), so \(\frac{1}{\sqrt{t^{2}+5}}\to 0\). Therefore the limit exists and
\[
\int_{2}^{\infty} \dfrac{x}{(x^{2}+5)^{\frac{3}{2}}}\,\mathrm{d}x=\lim_{t\to\infty}\left(\dfrac{1}{3}-\dfrac{1}{\sqrt{t^{2}+5}}\right)=\dfrac{1}{3}-0=\dfrac{1}{3}
\]

as required.
\end{djsSolution}

\newpage

% ---- Question 3 ----
\questionitem
% [Source: _QBT___Solns__Improper_Integrals_and_Mean_Values, Q3]

By first expressing
\[\dfrac{4}{x(x+1)(x+2)}\]
in partial fractions, show that the improper integral
\[\int_1^{\infty} \dfrac{4}{x(x+1)(x+2)}\,\mathrm{d}x\]
converges, and find its exact value in the form $\ln k$, where $k$ is a rational number to be determined. \marks{8} \\

\vspace*{10pt}

\begin{djsSolution}
Write the integrand in partial fractions:
\[
\dfrac{4}{x(x+1)(x+2)}=\dfrac{A}{x}+\dfrac{B}{x+1}+\dfrac{C}{x+2}
\]

Multiplying through by \(x(x+1)(x+2)\) gives
\[
4=A(x+1)(x+2)+Bx(x+2)+Cx(x+1)
\]

Substituting values of \(x\):
\begin{align*}
x=0:&\quad 4=2A\quad\Longrightarrow\quad A=2\\
x=-1:&\quad 4=-B\quad\Longrightarrow\quad B=-4\\
x=-2:&\quad 4=2C\quad\Longrightarrow\quad C=2
\end{align*}

So
\[
\dfrac{4}{x(x+1)(x+2)}=\dfrac{2}{x}-\dfrac{4}{x+1}+\dfrac{2}{x+2}
\]

The integral has an infinite upper limit, so replace it with \(t\) and consider the limit as \(t\to\infty\). For \(x\ge 1\), the expressions \(x\), \(x+1\) and \(x+2\) are all positive, so no modulus signs are needed in the logarithms. For \(t>1\):
\begin{align*}
\int_1^{t}\dfrac{4}{x(x+1)(x+2)}\,\mathrm{d}x
&=\int_1^{t}\left(\dfrac{2}{x}-\dfrac{4}{x+1}+\dfrac{2}{x+2}\right)\mathrm{d}x\\
&=\Big[2\ln x-4\ln(x+1)+2\ln(x+2)\Big]_1^{t}\\
&=\left[2\ln\left(\dfrac{x(x+2)}{(x+1)^2}\right)\right]_1^{t}\\
&=2\ln\left(\dfrac{t(t+2)}{(t+1)^2}\right)-2\ln\left(\dfrac{1\times3}{2^2}\right)\\
&=2\ln\left(\dfrac{t^2+2t}{t^2+2t+1}\right)-2\ln\left(\dfrac{3}{4}\right)
\end{align*}

Dividing the numerator and denominator by \(t^2\):
\[
\dfrac{t^2+2t}{t^2+2t+1}=\dfrac{1+\frac{2}{t}}{1+\frac{2}{t}+\frac{1}{t^2}}\to\dfrac{1+0}{1+0+0}=1\quad\text{as } t\to\infty
\]

Hence
\[
2\ln\left(\dfrac{t^2+2t}{t^2+2t+1}\right)\to 2\ln 1=0\quad\text{as } t\to\infty
\]

The limit exists and is finite, so the improper integral converges. Its value is
\begin{align*}
\int_1^{\infty}\dfrac{4}{x(x+1)(x+2)}\,\mathrm{d}x
&=\lim_{t\to\infty}\left[2\ln\left(\dfrac{t^2+2t}{t^2+2t+1}\right)-2\ln\left(\dfrac{3}{4}\right)\right]\\
&=0-2\ln\left(\dfrac{3}{4}\right)\\
&=2\ln\left(\dfrac{4}{3}\right)\\
&=\ln\left(\dfrac{16}{9}\right)
\end{align*}

This is in the form \(\ln k\) with \(k=\frac{16}{9}\).
\end{djsSolution}

\newpage

% ---- Question 4 ----
\questionitem
% [Source: _QBT___Solns__Improper_Integrals_and_Mean_Values, Q4]

Let
\[y = \dfrac{1}{1 - \sin x}\]
By first writing $y$ in terms of $\sec x$ and $\tan x$, find the exact mean value of $y$ with respect to $x$ over the interval $0 \le x \le \frac{1}{6}\pi$. \marks{6} \\

\vspace*{10pt}

\begin{djsSolution}
Multiply the numerator and denominator by \(1+\sin x\):
\begin{align*}
y&=\dfrac{1+\sin x}{(1-\sin x)(1+\sin x)}\\
&=\dfrac{1+\sin x}{1-\sin^2 x}\\
&=\dfrac{1+\sin x}{\cos^2 x}\\
&=\dfrac{1}{\cos^2 x}+\dfrac{1}{\cos x}\times\dfrac{\sin x}{\cos x}\\
&=\sec^2 x+\sec x\tan x
\end{align*}

The mean value of \(y\) over \(a\le x\le b\) is
\[
\dfrac{1}{b-a}\int_a^b y\,\mathrm{d}x
\]

Using \(\frac{\mathrm{d}}{\mathrm{d}x}(\tan x)=\sec^2 x\) and \(\frac{\mathrm{d}}{\mathrm{d}x}(\sec x)=\sec x\tan x\):
\begin{align*}
\int_0^{\frac{\pi}{6}}\left(\sec^2 x+\sec x\tan x\right)\mathrm{d}x&=\left[\tan x+\sec x\right]_0^{\frac{\pi}{6}}\\
&=\left(\dfrac{1}{\sqrt{3}}+\dfrac{2}{\sqrt{3}}\right)-(0+1)\\
&=\dfrac{3}{\sqrt{3}}-1\\
&=\sqrt{3}-1
\end{align*}

The interval has length \(\frac{\pi}{6}\), so the mean value is
\[
\dfrac{1}{\frac{\pi}{6}}\left(\sqrt{3}-1\right)=\dfrac{6\left(\sqrt{3}-1\right)}{\pi}
\]
\end{djsSolution}

\newpage

% ---- Question 5 ----
\questionitem
% [Source: _QBT___Solns__Improper_Integrals_and_Mean_Values, Q5]

The power drawn, $p$ watts, by an electric motor at time $t$ hours after it is switched on can be modelled by the equation
\[p = 40\tanh^2\left(\tfrac{1}{2}t\right), \qquad t \ge 0\] \\
Using an appropriate hyperbolic identity, find, according to the model, the \textbf{mean} power drawn by the motor over the first $2$ hours after it is switched on. Give your answer as an exact value in the form $\frac{a}{e^2+b}$, where $a$ and $b$ are integers. \marks{6} \\

\vspace*{10pt}

\begin{djsSolution}
The mean value of \(p\) over \(t_1\le t\le t_2\) is
\[
\dfrac{1}{t_2-t_1}\int_{t_1}^{t_2} p\,\mathrm{d}t
\]

So the mean value of \(p\) over \(0\le t\le 2\) is
\[
\frac{1}{2-0}\int_0^2 40\tanh^2\left(\tfrac{1}{2}t\right)\mathrm{d}t
\]

Dividing the identity \(\cosh^2 x-\sinh^2 x=1\) by \(\cosh^2 x\) gives
\[
1-\tanh^2 x=\operatorname{sech}^2 x\quad\Longrightarrow\quad \tanh^2 x=1-\operatorname{sech}^2 x
\]

So, with \(x=\frac{1}{2}t\), the mean value is
\[
20\int_0^2 \left(1-\operatorname{sech}^2\left(\tfrac{1}{2}t\right)\right)\mathrm{d}t
\]

Since \(\frac{\mathrm{d}}{\mathrm{d}t}\tanh\left(\frac{1}{2}t\right)=\frac{1}{2}\operatorname{sech}^2\left(\frac{1}{2}t\right)\), we have \(\int\operatorname{sech}^2\left(\frac{1}{2}t\right)\mathrm{d}t=2\tanh\left(\frac{1}{2}t\right)+c\). Hence the mean value is
\begin{align*}
20\left[t-2\tanh\left(\tfrac{1}{2}t\right)\right]_0^2&=20\left(2-2\tanh 1\right)-20(0-0)\\
&=40\left(1-\tanh 1\right)
\end{align*}

Now write \(\tanh 1\) in exponential form:
\[
\tanh 1=\frac{e-e^{-1}}{e+e^{-1}}=\frac{e^2-1}{e^2+1}
\]

so
\[
1-\tanh 1=\frac{(e^2+1)-(e^2-1)}{e^2+1}=\frac{2}{e^2+1}
\]

Therefore the mean value is
\[
40\times\frac{2}{e^2+1}=\frac{80}{e^2+1}
\]

The mean power drawn over the first \(2\) hours is \(\frac{80}{e^2+1}\,\mathrm{W}\), which is in the form \(\frac{a}{e^2+b}\) with \(a=80\), \(b=1\).
\end{djsSolution}

\newpage

% ---- Question 6 ----
\questionitem
% [Source: _QBT___Solns__Improper_Integrals_and_Mean_Values, Q6]

\begin{enumerate}[label=\textbf{(\alph*)}]
  \item Explain why
  \[\int_{e}^{\infty} \dfrac{\ln x}{x^{3}}\,\mathrm{d}x\]
  is an improper integral. \marks{1} \\
  \item Use integration by parts to find
  \[\int \dfrac{\ln x}{x^{3}}\,\mathrm{d}x\] \marks[-30pt]{3}
  \item Hence show that
  \[\int_{e}^{\infty} \dfrac{\ln x}{x^{3}}\,\mathrm{d}x\]
  converges, and find its exact value. You should show the limiting process used, and state clearly the behaviour of any limit you rely on. \marks{4} \\
\end{enumerate}

\vspace*{10pt}

\begin{djsSolution}
\begin{enumerate}[label=\textbf{(\alph*)}, start=1]
  \item The upper limit of integration is infinite, so the interval of integration is unbounded and the integral is improper.

  \item Integration by parts states
\[
\int u\frac{\mathrm{d}v}{\mathrm{d}x}\,\mathrm{d}x = uv - \int v\frac{\mathrm{d}u}{\mathrm{d}x}\,\mathrm{d}x
\]

Take \(u=\ln x\), \(\frac{\mathrm{d}v}{\mathrm{d}x}=x^{-3}\), so \(\frac{\mathrm{d}u}{\mathrm{d}x}=\frac{1}{x}\), \(v=-\frac{1}{2x^2}\). Then
\begin{align*}
\int \dfrac{\ln x}{x^{3}}\,\mathrm{d}x &= -\dfrac{\ln x}{2x^2} - \int \left(-\dfrac{1}{2x^2}\right)\dfrac{1}{x}\,\mathrm{d}x\\
&= -\dfrac{\ln x}{2x^2} + \dfrac12\int x^{-3}\,\mathrm{d}x\\
&= -\dfrac{\ln x}{2x^2} + \dfrac12\left(-\dfrac{1}{2x^2}\right) + c\\
&= -\dfrac{\ln x}{2x^2} - \dfrac{1}{4x^2} + c
\end{align*}

  \item Replace the infinite upper limit by \(t\), where \(t>e\), and let \(t\to\infty\):
\[
\int_{e}^{\infty} \dfrac{\ln x}{x^{3}}\,\mathrm{d}x = \lim_{t\to\infty}\int_{e}^{t} \dfrac{\ln x}{x^{3}}\,\mathrm{d}x
\]

Using part (b), and \(\ln e = 1\):
\begin{align*}
\int_{e}^{t} \dfrac{\ln x}{x^{3}}\,\mathrm{d}x &= \left[-\dfrac{\ln x}{2x^2} - \dfrac{1}{4x^2}\right]_{e}^{t}\\
&= \left(-\dfrac{\ln t}{2t^2} - \dfrac{1}{4t^2}\right) - \left(-\dfrac{1}{2e^2} - \dfrac{1}{4e^2}\right)\\
&= \dfrac{3}{4e^2} - \dfrac{\ln t}{2t^2} - \dfrac{1}{4t^2}
\end{align*}

As \(t\to\infty\), \(\frac{\ln t}{t^2}\to 0\), since a positive power of \(t\) grows faster than \(\ln t\). Also \(\frac{1}{t^2}\to 0\). Therefore
\[
\lim_{t\to\infty}\int_{e}^{t} \dfrac{\ln x}{x^{3}}\,\mathrm{d}x = \dfrac{3}{4e^2} - 0 - 0 = \dfrac{3}{4e^2}
\]

The limit exists and is finite, so the integral converges, with
\[
\int_{e}^{\infty} \dfrac{\ln x}{x^{3}}\,\mathrm{d}x = \dfrac{3}{4e^2}
\]

\end{enumerate}
\end{djsSolution}

\newpage

% ---- Question 7 ----
\questionitem
% [Source: _QBT___Solns__Improper_Integrals_and_Mean_Values, Q7]

\begin{enumerate}[label=\textbf{(\alph*)}]
  \item Express $9x^2 - 18x + 25$ in the form $a(x+p)^2 + q$, where $a$, $p$ and $q$ are integers. \marks{2} \\
  \item Hence use a standard integral from the formula book to find
  \[\int \dfrac{1}{\sqrt{9x^2 - 18x + 25}}\,\mathrm{d}x\] \marks[-30pt]{2}
  \item Determine the exact mean value of the function
  \[\mathrm{g}(x) = \dfrac{1}{\sqrt{9x^2 - 18x + 25}}, \qquad 0 \le x \le 2\]
  giving your answer in the form $A \ln B$, where $A$ and $B$ are rational constants with $B$ an integer. \marks{3} \\
\end{enumerate}

\vspace*{10pt}

\begin{djsSolution}
\begin{enumerate}[label=\textbf{(\alph*)}, start=1]
  \item Take out the factor of \(9\) from the \(x\) terms and complete the square:
\begin{align*}
9x^2-18x+25&=9\left(x^2-2x\right)+25\\
&=9\left[(x-1)^2-1\right]+25\\
&=9(x-1)^2+16
\end{align*}

This is in the form \(a(x+p)^2+q\) with \(a=9\), \(p=-1\), \(q=16\).

  \item Using part (a), factor out the \(9\) inside the square root:
\begin{align*}
\int \dfrac{1}{\sqrt{9x^2-18x+25}}\,\mathrm{d}x&=\int \dfrac{1}{\sqrt{9(x-1)^2+16}}\,\mathrm{d}x\\
&=\dfrac{1}{3}\int \dfrac{1}{\sqrt{(x-1)^2+\left(\frac{4}{3}\right)^2}}\,\mathrm{d}x
\end{align*}

The standard result from the formula book, with its constant written as \(k\) since \(a\) is already in use, is
\[
\int \dfrac{1}{\sqrt{x^2+k^2}}\,\mathrm{d}x=\operatorname{arsinh}\left(\dfrac{x}{k}\right)+c
\]

Applying this with \(k=\frac{4}{3}\), and with \(x\) replaced by \(x-1\):
\begin{align*}
\int \dfrac{1}{\sqrt{9x^2-18x+25}}\,\mathrm{d}x&=\dfrac{1}{3}\operatorname{arsinh}\left(\dfrac{x-1}{\frac{4}{3}}\right)+c\\
&=\dfrac{1}{3}\operatorname{arsinh}\left(\dfrac{3(x-1)}{4}\right)+c
\end{align*}

  \item The mean value of \(\mathrm{g}(x)\) over \(x_1\le x\le x_2\) is
\[
\dfrac{1}{x_2-x_1}\int_{x_1}^{x_2} \mathrm{g}(x)\,\mathrm{d}x
\]

So the mean value of \(\mathrm{g}\) over \(0\le x\le 2\) is
\begin{align*}
\dfrac{1}{2-0}\int_0^2 \mathrm{g}(x)\,\mathrm{d}x&=\dfrac{1}{2}\times\dfrac{1}{3}\left[\operatorname{arsinh}\left(\dfrac{3(x-1)}{4}\right)\right]_0^2\\
&=\dfrac{1}{6}\left(\operatorname{arsinh}\left(\dfrac{3}{4}\right)-\operatorname{arsinh}\left(-\dfrac{3}{4}\right)\right)
\end{align*}

The formula book gives
\[
\operatorname{arsinh} t=\ln\left(t+\sqrt{t^2+1}\right)
\]

Since \(\sqrt{\frac{9}{16}+1}=\frac{5}{4}\):
\begin{align*}
\operatorname{arsinh}\left(\dfrac{3}{4}\right)&=\ln\left(\dfrac{3}{4}+\dfrac{5}{4}\right)=\ln 2\\
\operatorname{arsinh}\left(-\dfrac{3}{4}\right)&=\ln\left(-\dfrac{3}{4}+\dfrac{5}{4}\right)=\ln\dfrac{1}{2}=-\ln 2
\end{align*}

Therefore the mean value is
\[
\dfrac{1}{6}\left(\ln 2-(-\ln 2)\right)=\dfrac{1}{6}\times 2\ln 2=\dfrac{1}{3}\ln 2
\]

This is in the form \(A\ln B\) with \(A=\frac{1}{3}\), \(B=2\).
\end{enumerate}
\end{djsSolution}

\newpage

% ---- Question 8 ----
\questionitem
% [Source: _QBT___Solns__Improper_Integrals_and_Mean_Values, Q8]

Let $y = \ln(1 + x)$. \\
\begin{enumerate}[label=\textbf{(\alph*)}]
  \item Show that the mean value of $y$ with respect to $x$ over the interval $0 \le x \le e - 1$ is
  \[\dfrac{1}{e - 1}\] \marks[-30pt]{4}
  \item Find, in exact form, the mean value of $x$ with respect to $y$ over the interval $0 \le y \le 1$. \marks{3} \\
\end{enumerate}

\vspace*{10pt}

\begin{djsSolution}
\begin{enumerate}[label=\textbf{(\alph*)}, start=1]
  \item The mean value of \(f(x)\) over \(a\le x\le b\) is
\[
\dfrac{1}{b-a}\int_a^b f(x)\,\mathrm{d}x
\]

So the integral \(\int_0^{e-1}\ln(1+x)\,\mathrm{d}x\) is needed. Integration by parts states
\[
\displaystyle\int u\frac{\mathrm{d}v}{\mathrm{d}x}\,\mathrm{d}x = uv - \int v\frac{\mathrm{d}u}{\mathrm{d}x}\,\mathrm{d}x
\]

Take \(u=\ln(1+x)\), \(\frac{\mathrm{d}v}{\mathrm{d}x}=1\), so \(\frac{\mathrm{d}u}{\mathrm{d}x}=\frac{1}{1+x}\), \(v=x\). Note that \(1+x>0\) on the interval. Then
\begin{align*}
\int \ln(1+x)\,\mathrm{d}x &= x\ln(1+x)-\int\frac{x}{1+x}\,\mathrm{d}x\\
&= x\ln(1+x)-\int\left(1-\frac{1}{1+x}\right)\mathrm{d}x\\
&= x\ln(1+x)-x+\ln(1+x)+c
\end{align*}

Evaluating between the limits, using \(\ln e=1\) and \(\ln 1=0\):
\begin{align*}
\int_0^{e-1}\ln(1+x)\,\mathrm{d}x &= \Big[(1+x)\ln(1+x)-x\Big]_0^{e-1}\\
&= \big(e\ln e-(e-1)\big)-\big(1\times\ln 1-0\big)\\
&= e-e+1\\
&= 1
\end{align*}

Hence the mean value is
\[
\frac{1}{(e-1)-0}\times 1=\frac{1}{e-1}
\]

as required.

  \item Rearranging \(y=\ln(1+x)\):
\[
1+x=e^{y}\quad\Longrightarrow\quad x=e^{y}-1
\]

The mean value of \(x\) with respect to \(y\) over \(0\le y\le 1\) is
\begin{align*}
\frac{1}{1-0}\int_0^1\left(e^{y}-1\right)\mathrm{d}y &= \Big[e^{y}-y\Big]_0^1\\
&= (e-1)-(1-0)\\
&= e-2
\end{align*}

\end{enumerate}
\end{djsSolution}

\newpage

% ---- Question 9 ----
\questionitem
% [Source: _QBT___Solns__Improper_Integrals_and_Mean_Values, Q9]

In an electrical circuit, the current $I$ amps at time $t$ seconds is modelled by
\[ I = a\sin nt + b\sin 2nt \]
where $a$, $b$ and $n$ are positive constants. The RMS value of the current, in amps, is defined to be the square root of the mean value of $I^2$ over the interval $0 \le t \le \frac{2\pi}{n}$. \\

Show that the RMS value of the current is
\[\sqrt{\dfrac{a^2+b^2}{2}}\]
amps. \marks{7} \\

\vspace*{10pt}

\begin{djsSolution}
The mean value of \(f(t)\) over \(p\le t\le q\) is
\[
\dfrac{1}{q-p}\int_p^q f(t)\,\mathrm{d}t
\]

So the mean value of \(I^2\) over \(0 \le t \le \frac{2\pi}{n}\) is
\[
\frac{n}{2\pi}\int_0^{2\pi/n} I^2\,\mathrm{d}t
\]

Expanding \(I^2\):
\[
I^2 = a^2\sin^2 nt + 2ab\sin nt\sin 2nt + b^2\sin^2 2nt
\]

Using the double angle formula \(\cos 2x = 1-2\sin^2 x\), we have \(\sin^2 x = \frac{1}{2}(1-\cos 2x)\). Hence
\begin{align*}
\int_0^{2\pi/n} \sin^2 nt\,\mathrm{d}t &= \int_0^{2\pi/n} \frac{1}{2}\left(1-\cos 2nt\right)\mathrm{d}t\\
&= \left[\frac{t}{2} - \frac{\sin 2nt}{4n}\right]_0^{2\pi/n}\\
&= \left(\frac{\pi}{n} - \frac{\sin 4\pi}{4n}\right) - 0\\
&= \frac{\pi}{n}
\end{align*}

Similarly
\begin{align*}
\int_0^{2\pi/n} \sin^2 2nt\,\mathrm{d}t &= \int_0^{2\pi/n} \frac{1}{2}\left(1-\cos 4nt\right)\mathrm{d}t\\
&= \left[\frac{t}{2} - \frac{\sin 4nt}{8n}\right]_0^{2\pi/n}\\
&= \left(\frac{\pi}{n} - \frac{\sin 8\pi}{8n}\right) - 0\\
&= \frac{\pi}{n}
\end{align*}

For the cross term, use the double angle formula \(\sin 2nt = 2\sin nt\cos nt\):
\begin{align*}
\int_0^{2\pi/n} \sin nt\sin 2nt\,\mathrm{d}t &= \int_0^{2\pi/n} 2\sin^2 nt\cos nt\,\mathrm{d}t\\
&= \left[\frac{2\sin^3 nt}{3n}\right]_0^{2\pi/n}\\
&= \frac{2\sin^3 2\pi}{3n} - 0\\
&= 0
\end{align*}

Combining these results:
\begin{align*}
\int_0^{2\pi/n} I^2\,\mathrm{d}t &= a^2\times\frac{\pi}{n} + 2ab\times 0 + b^2\times\frac{\pi}{n}\\
&= \frac{\left(a^2+b^2\right)\pi}{n}
\end{align*}

So the mean value of \(I^2\) is
\[
\frac{n}{2\pi}\times\frac{\left(a^2+b^2\right)\pi}{n} = \frac{a^2+b^2}{2}
\]

Taking the square root, the RMS value of the current is
\[
\sqrt{\dfrac{a^2+b^2}{2}}
\]

amps, as required.
\end{djsSolution}

\newpage

% ---- Question 10 ----
\questionitem
% [Source: _QBT___Solns__Improper_Integrals_and_Mean_Values, Q10]

Prove that
\[\int_{1}^{\infty} \dfrac{x+3}{x^{2}(x+1)}\,\mathrm{d}x = 3 - \ln 4\]
Show the limiting process clearly. \marks{8} \\

\vspace*{10pt}

\begin{djsSolution}
The integral is improper, so it is defined as
\[
\int_{1}^{\infty} \dfrac{x+3}{x^{2}(x+1)}\,\mathrm{d}x=\lim_{t\to\infty}\int_{1}^{t} \dfrac{x+3}{x^{2}(x+1)}\,\mathrm{d}x
\]

First write the integrand in partial fractions:
\[
\dfrac{x+3}{x^{2}(x+1)}=\dfrac{A}{x}+\dfrac{B}{x^2}+\dfrac{C}{x+1}\quad\Longrightarrow\quad x+3=Ax(x+1)+B(x+1)+Cx^2
\]

Setting \(x=0\) gives \(B=3\). Setting \(x=-1\) gives \(2=C\). Equating coefficients of \(x^2\) gives \(A+C=0\), so \(A=-2\). Hence
\[
\dfrac{x+3}{x^{2}(x+1)}=-\dfrac{2}{x}+\dfrac{3}{x^2}+\dfrac{2}{x+1}
\]

Since \(x\ge 1\), both \(x\) and \(x+1\) are positive, so no modulus signs are needed in the logarithms. For \(t>1\):
\begin{align*}
\int_{1}^{t} \left(-\dfrac{2}{x}+\dfrac{3}{x^2}+\dfrac{2}{x+1}\right)\mathrm{d}x
&=\left[-2\ln x-\dfrac{3}{x}+2\ln(x+1)\right]_{1}^{t}\\
&=\left[2\ln\left(\dfrac{x+1}{x}\right)-\dfrac{3}{x}\right]_{1}^{t}\\
&=\left(2\ln\left(1+\dfrac{1}{t}\right)-\dfrac{3}{t}\right)-\left(2\ln 2-3\right)
\end{align*}

The logarithms are combined before taking the limit, since \(\ln t\) and \(\ln(t+1)\) separately both tend to infinity.

As \(t\to\infty\), \(\frac{1}{t}\to 0\), so \(1+\frac{1}{t}\to 1\) and \(\ln\left(1+\frac{1}{t}\right)\to\ln 1=0\). Also \(\frac{3}{t}\to 0\). Therefore
\begin{align*}
\int_{1}^{\infty} \dfrac{x+3}{x^{2}(x+1)}\,\mathrm{d}x
&=\lim_{t\to\infty}\left(2\ln\left(1+\dfrac{1}{t}\right)-\dfrac{3}{t}-2\ln 2+3\right)\\
&=0-0-2\ln 2+3\\
&=3-\ln 4
\end{align*}

as required.
\end{djsSolution}

\newpage

% ---- Question 11 ----
\questionitem
% [Source: _QBT___Solns__Improper_Integrals_and_Mean_Values, Q11]

\begin{enumerate}[label=\textbf{(\alph*)}]
  \item Explain why
  \[\int_{0}^{\infty} \dfrac{10}{(x+2)(3x+1)}\,\mathrm{d}x\]
  is an improper integral. \marks{1} \\
  \item Show that
  \[\int_{0}^{\infty} \dfrac{10}{(x+2)(3x+1)}\,\mathrm{d}x = a\ln b\]
  where $a$ is an integer and $b$ is a rational number, both to be determined. You must show clearly how you deal with the infinite limit. \marks{6} \\
\end{enumerate}

\vspace*{10pt}

\begin{djsSolution}
\begin{enumerate}[label=\textbf{(\alph*)}, start=1]
  \item The upper limit of integration is infinite, so the interval of integration is unbounded and the integral is improper.
  \item First express the integrand in partial fractions:
\[
\dfrac{10}{(x+2)(3x+1)}=\dfrac{A}{x+2}+\dfrac{B}{3x+1}
\]

so
\[
10=A(3x+1)+B(x+2)
\]

Setting \(x=-2\) gives \(10=-5A\), so \(A=-2\). Setting \(x=-\frac{1}{3}\) gives \(10=\frac{5}{3}B\), so \(B=6\).

For \(x\ge 0\), both \(x+2\) and \(3x+1\) are positive, so
\begin{align*}
\int \dfrac{10}{(x+2)(3x+1)}\,\mathrm{d}x &= \int\left(-\dfrac{2}{x+2}+\dfrac{6}{3x+1}\right)\mathrm{d}x\\
&=-2\ln(x+2)+2\ln(3x+1)+c\\
&=2\ln\left(\dfrac{3x+1}{x+2}\right)+c
\end{align*}

Replace the infinite limit by \(t\), where \(t>0\):
\begin{align*}
\int_0^t \dfrac{10}{(x+2)(3x+1)}\,\mathrm{d}x &= \left[2\ln\left(\dfrac{3x+1}{x+2}\right)\right]_0^t\\
&=2\ln\left(\dfrac{3t+1}{t+2}\right)-2\ln\left(\dfrac{1}{2}\right)\\
&=2\ln\left(\dfrac{3t+1}{t+2}\right)+2\ln 2
\end{align*}

Dividing the numerator and denominator by \(t\):
\[
\dfrac{3t+1}{t+2}=\dfrac{3+\frac{1}{t}}{1+\frac{2}{t}}\to\dfrac{3+0}{1+0}=3\quad\text{as } t\to\infty
\]

so \(2\ln\left(\frac{3t+1}{t+2}\right)\to 2\ln 3\). Therefore
\begin{align*}
\int_0^\infty \dfrac{10}{(x+2)(3x+1)}\,\mathrm{d}x &= \lim_{t\to\infty}\left(2\ln\left(\dfrac{3t+1}{t+2}\right)+2\ln 2\right)\\
&=2\ln 3+2\ln 2\\
&=2\ln 6
\end{align*}

This is in the form \(a\ln b\) with \(a=2\), \(b=6\).
\end{enumerate}
\end{djsSolution}

\newpage

% ---- Question 12 ----
\questionitem
% [Source: _QBT___Solns__Improper_Integrals_and_Mean_Values, Q12]

\begin{enumerate}[label=\textbf{(\alph*)}]
  \item Explain why
  \[\int_0^3 \dfrac{x}{\sqrt[3]{x^2-1}}\,\mathrm{d}x\]
  is an improper integral. \marks{1} \\
  \item Use an appropriate limit argument to show that this integral converges, and find its exact value. \marks{5} \\
\end{enumerate}

\vspace*{10pt}

\begin{djsSolution}
\begin{enumerate}[label=\textbf{(\alph*)}, start=1]
  \item When \(x=1\), \(x^2-1=0\), so the denominator \(\sqrt[3]{x^2-1}\) is zero and the integrand is undefined (it is unbounded near \(x=1\)). Since \(x=1\) lies inside the interval of integration \([0,3]\), the integral is improper.

  \item First find an antiderivative. Write the integrand as \(x\left(x^2-1\right)^{-\frac13}\), taking real cube roots, which exist for negative numbers too: for example \(\sqrt[3]{-1}=-1\). By the chain rule, for \(x\neq 1\):
\[
\frac{\mathrm{d}}{\mathrm{d}x}\left(\sqrt[3]{x^2-1}\right)^2=\frac{\mathrm{d}}{\mathrm{d}x}\left(x^2-1\right)^{\frac23}=\dfrac23\left(x^2-1\right)^{-\frac13}\times 2x=\dfrac43\,\dfrac{x}{\sqrt[3]{x^2-1}}
\]

Hence
\[
\int \dfrac{x}{\sqrt[3]{x^2-1}}\,\mathrm{d}x=\dfrac34\left(\sqrt[3]{x^2-1}\right)^2+c
\]

Split the integral at the point where the integrand is undefined:
\[
\int_0^3 \dfrac{x}{\sqrt[3]{x^2-1}}\,\mathrm{d}x=\int_0^1 \dfrac{x}{\sqrt[3]{x^2-1}}\,\mathrm{d}x+\int_1^3 \dfrac{x}{\sqrt[3]{x^2-1}}\,\mathrm{d}x
\]

The original integral converges if and only if both of these converge.

For the first integral, with \(0<t<1\):
\begin{align*}
\int_0^1 \dfrac{x}{\sqrt[3]{x^2-1}}\,\mathrm{d}x&=\lim_{t\to 1^-}\left[\dfrac34\left(\sqrt[3]{x^2-1}\right)^2\right]_0^t\\
&=\lim_{t\to 1^-}\left(\dfrac34\left(\sqrt[3]{t^2-1}\right)^2-\dfrac34\left(\sqrt[3]{-1}\right)^2\right)\\
&=0-\dfrac34\times(-1)^2\\
&=-\dfrac34
\end{align*}

since \(t^2-1\to 0\) as \(t\to1^-\). This is finite, so the first integral converges.

For the second integral, with \(1<s<3\):
\begin{align*}
\int_1^3 \dfrac{x}{\sqrt[3]{x^2-1}}\,\mathrm{d}x&=\lim_{s\to 1^+}\left[\dfrac34\left(\sqrt[3]{x^2-1}\right)^2\right]_s^3\\
&=\lim_{s\to 1^+}\left(\dfrac34\left(\sqrt[3]{8}\right)^2-\dfrac34\left(\sqrt[3]{s^2-1}\right)^2\right)\\
&=\dfrac34\times 4-0\\
&=3
\end{align*}

This is also finite, so the second integral converges.

Both limits exist and are finite, so the integral converges, and
\[
\int_0^3 \dfrac{x}{\sqrt[3]{x^2-1}}\,\mathrm{d}x=-\dfrac34+3=\dfrac94
\]

\end{enumerate}
\end{djsSolution}

\newpage

% ---- Question 13 ----
\questionitem
% [Source: _QBT___Solns__Improper_Integrals_and_Mean_Values, Q13]

\begin{enumerate}[label=\textbf{(\alph*)}]
  \item Explain why
  \[\int_0^{e} \dfrac{(\ln x)^2}{\sqrt{x}}\,\mathrm{d}x\]
  is an improper integral. \marks{1} \\
  \item By using integration by parts twice, show that
  \[\int x^{-\frac{1}{2}} (\ln x)^2\,\mathrm{d}x = 2\sqrt{x}\left((\ln x)^2 - 4\ln x + 8\right) + c\] \marks[-30pt]{4}
  \item You may assume that, for any constant $k > 0$ and any positive integer $n$, $x^k(\ln x)^n \to 0$ as $x \to 0^{+}$. \\

  Show that
  \[\int_0^{e} \dfrac{(\ln x)^2}{\sqrt{x}}\,\mathrm{d}x\]
  exists and find its exact value. \marks{4} \\
\end{enumerate}

\vspace*{10pt}

\begin{djsSolution}
\begin{enumerate}[label=\textbf{(\alph*)}, start=1]
  \item The integrand is not defined at the lower limit \(x=0\). As \(x\to0^{+}\), \((\ln x)^2\to\infty\) and \(\sqrt{x}\to0\), so \(\frac{(\ln x)^2}{\sqrt{x}}\to\infty\). The integrand is unbounded on the interval of integration, so the integral is improper.

  \item Integration by parts states
\[
\int u\frac{\mathrm{d}v}{\mathrm{d}x}\,\mathrm{d}x = uv-\int v\frac{\mathrm{d}u}{\mathrm{d}x}\,\mathrm{d}x
\]

First take \(u=(\ln x)^2\), \(\frac{\mathrm{d}v}{\mathrm{d}x}=x^{-\frac12}\), so \(\frac{\mathrm{d}u}{\mathrm{d}x}=\frac{2\ln x}{x}\), \(v=2x^{\frac12}\). Then
\begin{align*}
\int x^{-\frac12}(\ln x)^2\,\mathrm{d}x &= 2x^{\frac12}(\ln x)^2-\int 2x^{\frac12}\times\frac{2\ln x}{x}\,\mathrm{d}x\\
&= 2\sqrt{x}(\ln x)^2-4\int x^{-\frac12}\ln x\,\mathrm{d}x
\end{align*}

For the remaining integral take \(u=\ln x\), \(\frac{\mathrm{d}v}{\mathrm{d}x}=x^{-\frac12}\), so \(\frac{\mathrm{d}u}{\mathrm{d}x}=\frac{1}{x}\), \(v=2x^{\frac12}\). Then
\begin{align*}
\int x^{-\frac12}\ln x\,\mathrm{d}x &= 2x^{\frac12}\ln x-\int 2x^{\frac12}\times\frac{1}{x}\,\mathrm{d}x\\
&= 2\sqrt{x}\ln x-2\int x^{-\frac12}\,\mathrm{d}x\\
&= 2\sqrt{x}\ln x-4\sqrt{x}
\end{align*}

Combining the two results:
\begin{align*}
\int x^{-\frac12}(\ln x)^2\,\mathrm{d}x &= 2\sqrt{x}(\ln x)^2-4\left(2\sqrt{x}\ln x-4\sqrt{x}\right)+c\\
&= 2\sqrt{x}(\ln x)^2-8\sqrt{x}\ln x+16\sqrt{x}+c\\
&= 2\sqrt{x}\left((\ln x)^2-4\ln x+8\right)+c
\end{align*}

as required.

  \item Since the integrand is unbounded at \(x=0\), replace the lower limit by \(t\), where \(0<t<e\), and let \(t\to0^{+}\):
\[
\int_0^{e}\dfrac{(\ln x)^2}{\sqrt{x}}\,\mathrm{d}x=\lim_{t\to0^{+}}\int_t^{e}x^{-\frac12}(\ln x)^2\,\mathrm{d}x
\]

Using part (b):
\begin{align*}
\int_t^{e}x^{-\frac12}(\ln x)^2\,\mathrm{d}x &= \left[2\sqrt{x}\left((\ln x)^2-4\ln x+8\right)\right]_t^{e}\\
&= 2\sqrt{e}\,(1-4+8)-2\sqrt{t}\left((\ln t)^2-4\ln t+8\right)\\
&= 10\sqrt{e}-2t^{\frac12}(\ln t)^2+8t^{\frac12}\ln t-16t^{\frac12}
\end{align*}

As \(t\to0^{+}\), using the given result with \(k=\frac12\) and \(n=2\), \(t^{\frac12}(\ln t)^2\to0\); with \(k=\frac12\) and \(n=1\), \(t^{\frac12}\ln t\to0\); and \(16t^{\frac12}\to0\). Hence
\[
\lim_{t\to0^{+}}\int_t^{e}x^{-\frac12}(\ln x)^2\,\mathrm{d}x=10\sqrt{e}-0+0-0=10\sqrt{e}
\]

The limit is finite, so the integral exists, and
\[
\int_0^{e}\dfrac{(\ln x)^2}{\sqrt{x}}\,\mathrm{d}x=10\sqrt{e}
\]

\end{enumerate}
\end{djsSolution}

\newpage

% ---- Question 14 ----
\questionitem
% [Source: _QBT___Solns__Improper_Integrals_and_Mean_Values, Q14]

\begin{enumerate}[label=\textbf{(\alph*)}]
  \item Find
  \[\int_1^a \dfrac{(\ln x)^2}{x^2}\,\mathrm{d}x\]
  where $a > 1$, giving your answer in terms of $a$. \marks{6} \\
  \item Write down the value of
  \[\lim_{a \to \infty} \dfrac{(\ln a)^k}{a}\]
  where $k$ is a positive constant. \marks{1} \\
  \item Hence find
  \[\int_1^{\infty} \dfrac{(\ln x)^2}{x^2}\,\mathrm{d}x\] \marks[-30pt]{2}
\end{enumerate}

\vspace*{10pt}

\begin{djsSolution}
\begin{enumerate}[label=\textbf{(\alph*)}, start=1]
  \item Integration by parts states
\[
\int u\frac{\mathrm{d}v}{\mathrm{d}x}\,\mathrm{d}x = uv - \int v\frac{\mathrm{d}u}{\mathrm{d}x}\,\mathrm{d}x
\]

Take \(u=(\ln x)^2\), \(\frac{\mathrm{d}v}{\mathrm{d}x}=x^{-2}\), so \(\frac{\mathrm{d}u}{\mathrm{d}x}=\frac{2\ln x}{x}\), \(v=-\frac{1}{x}\). Then
\begin{align*}
\int \frac{(\ln x)^2}{x^2}\,\mathrm{d}x &= -\frac{(\ln x)^2}{x} - \int \left(-\frac{1}{x}\right)\frac{2\ln x}{x}\,\mathrm{d}x\\
&= -\frac{(\ln x)^2}{x} + 2\int \frac{\ln x}{x^2}\,\mathrm{d}x
\end{align*}

For the remaining integral, apply integration by parts again with \(u=\ln x\), \(\frac{\mathrm{d}v}{\mathrm{d}x}=x^{-2}\), so \(\frac{\mathrm{d}u}{\mathrm{d}x}=\frac{1}{x}\), \(v=-\frac{1}{x}\):
\begin{align*}
\int \frac{\ln x}{x^2}\,\mathrm{d}x &= -\frac{\ln x}{x} - \int \left(-\frac{1}{x}\right)\frac{1}{x}\,\mathrm{d}x\\
&= -\frac{\ln x}{x} + \int x^{-2}\,\mathrm{d}x\\
&= -\frac{\ln x}{x} - \frac{1}{x} + c
\end{align*}

Combining these:
\[
\int \frac{(\ln x)^2}{x^2}\,\mathrm{d}x = -\frac{(\ln x)^2}{x} - \frac{2\ln x}{x} - \frac{2}{x} + c
\]

Evaluating between the limits, using \(\ln 1=0\):
\begin{align*}
\int_1^a \frac{(\ln x)^2}{x^2}\,\mathrm{d}x &= \left[-\frac{(\ln x)^2}{x} - \frac{2\ln x}{x} - \frac{2}{x}\right]_1^a\\
&= \left(-\frac{(\ln a)^2}{a} - \frac{2\ln a}{a} - \frac{2}{a}\right) - (0 - 0 - 2)\\
&= 2 - \frac{(\ln a)^2}{a} - \frac{2\ln a}{a} - \frac{2}{a}
\end{align*}

  \item Any power of \(\ln a\) grows more slowly than \(a\), so
\[
\lim_{a \to \infty} \frac{(\ln a)^k}{a} = 0
\]

  \item The improper integral is defined as a limit, so using part (a):
\[
\int_1^{\infty} \frac{(\ln x)^2}{x^2}\,\mathrm{d}x = \lim_{a\to\infty}\left(2 - \frac{(\ln a)^2}{a} - \frac{2\ln a}{a} - \frac{2}{a}\right)
\]

By part (b) with \(k=2\) and with \(k=1\), \(\frac{(\ln a)^2}{a}\to 0\) and \(\frac{\ln a}{a}\to 0\) as \(a\to\infty\); also \(\frac{2}{a}\to 0\). Therefore
\[
\int_1^{\infty} \frac{(\ln x)^2}{x^2}\,\mathrm{d}x = 2 - 0 - 0 - 0 = 2
\]

\end{enumerate}
\end{djsSolution}

\newpage

% ---- Question 15 ----
\questionitem
% [Source: _QBT___Solns__Improper_Integrals_and_Mean_Values, Q15]

Show that
\[\int_0^{\infty} \dfrac{5x-3}{(x^2+3)(x+1)}\,\mathrm{d}x = p\sqrt{3}\,\pi + \ln q\]
where $p$ and $q$ are rational numbers to be found. \marks{8} \\

\vspace*{10pt}

\begin{djsSolution}
First express the integrand in partial fractions:
\[
\dfrac{5x-3}{(x^2+3)(x+1)}=\dfrac{A}{x+1}+\dfrac{Bx+C}{x^2+3}
\]

Multiplying through by \((x^2+3)(x+1)\) gives
\[
5x-3=A(x^2+3)+(Bx+C)(x+1)
\]

Setting \(x=-1\):
\[
-8=4A\quad\Longrightarrow\quad A=-2
\]

Equating coefficients of \(x^2\) and the constant terms:
\begin{align*}
x^2:&\quad 0=A+B\quad\Longrightarrow\quad B=2\\
\text{constant}:&\quad -3=3A+C\quad\Longrightarrow\quad C=3
\end{align*}

So
\[
\dfrac{5x-3}{(x^2+3)(x+1)}=-\dfrac{2}{x+1}+\dfrac{2x}{x^2+3}+\dfrac{3}{x^2+3}
\]

For the last term, the standard result is
\[
\int\dfrac{1}{x^2+a^2}\,\mathrm{d}x=\dfrac{1}{a}\arctan\left(\dfrac{x}{a}\right)+c
\]

with \(a=\sqrt{3}\). Since \(x+1>0\) for \(x\ge 0\), the logarithm needs no modulus, so
\begin{align*}
\int\dfrac{5x-3}{(x^2+3)(x+1)}\,\mathrm{d}x&=-2\ln(x+1)+\ln(x^2+3)+\dfrac{3}{\sqrt{3}}\arctan\left(\dfrac{x}{\sqrt{3}}\right)+c\\
&=\ln\left(\dfrac{x^2+3}{(x+1)^2}\right)+\sqrt{3}\arctan\left(\dfrac{x}{\sqrt{3}}\right)+c
\end{align*}

The integral is improper, so write it as a limit:
\begin{align*}
\int_0^{\infty}\dfrac{5x-3}{(x^2+3)(x+1)}\,\mathrm{d}x&=\lim_{t\to\infty}\left[\ln\left(\dfrac{x^2+3}{(x+1)^2}\right)+\sqrt{3}\arctan\left(\dfrac{x}{\sqrt{3}}\right)\right]_0^{t}\\
&=\lim_{t\to\infty}\left[\ln\left(\dfrac{t^2+3}{(t+1)^2}\right)+\sqrt{3}\arctan\left(\dfrac{t}{\sqrt{3}}\right)\right]-\left(\ln 3+\sqrt{3}\arctan 0\right)
\end{align*}

Dividing the numerator and denominator by \(t^2\):
\[
\dfrac{t^2+3}{(t+1)^2}=\dfrac{1+\frac{3}{t^2}}{\left(1+\frac{1}{t}\right)^2}\to 1\quad\text{as } t\to\infty
\]

so \(\ln\left(\frac{t^2+3}{(t+1)^2}\right)\to\ln 1=0\). Also \(\arctan\left(\frac{t}{\sqrt{3}}\right)\to\frac{\pi}{2}\), and \(\arctan 0=0\). Therefore
\begin{align*}
\int_0^{\infty}\dfrac{5x-3}{(x^2+3)(x+1)}\,\mathrm{d}x&=0+\dfrac{\sqrt{3}\,\pi}{2}-\ln 3\\
&=\dfrac{1}{2}\sqrt{3}\,\pi+\ln\dfrac{1}{3}
\end{align*}

This is in the form \(p\sqrt{3}\,\pi+\ln q\) with \(p=\frac{1}{2}\) and \(q=\frac{1}{3}\).
\end{djsSolution}

\newpage

% ---- Question 16 ----
\questionitem
% [Source: _QBT___Solns__Improper_Integrals_and_Mean_Values, Q16]

\begin{enumerate}[label=\textbf{(\alph*)}]
  \item Explain why
  \[\int_{0}^{1} (\ln y)^2\,\mathrm{d}y\]
  is an improper integral. \marks{1} \\
  \item
  \begin{enumerate}[label=(\roman*)]
    \item Show that the substitution $y = e^{-x}$ transforms
    \[\int_{0}^{1} (\ln y)^2\,\mathrm{d}y \quad \text{into} \quad \int_{0}^{\infty} x^2 e^{-x}\,\mathrm{d}x\]
    explaining clearly how the limits of integration are obtained. \marks{3} \\
    \item Evaluate
    \[\int_{0}^{\infty} x^2 e^{-x}\,\mathrm{d}x\]
    showing the limiting process used. \marks{6} \\
    \item Hence write down the value of
    \[\int_{0}^{1} (\ln y)^2\,\mathrm{d}y\] \marks[-30pt]{1}
  \end{enumerate}
\end{enumerate}

\vspace*{10pt}

\begin{djsSolution}
\begin{enumerate}[label=\textbf{(\alph*)}, start=1]
  \item As \(y \to 0^{+}\), \(\ln y \to -\infty\), so \((\ln y)^2 \to \infty\). The integrand is therefore unbounded (and undefined) at the lower limit \(y=0\), so the integral is improper.

  \item
  \begin{enumerate}[label=(\roman*), start=1]
    \item With \(y=e^{-x}\), we have \(\ln y=-x\), so \((\ln y)^2=x^2\). Also
\[
\frac{\mathrm{d}y}{\mathrm{d}x}=-e^{-x}\quad\Longrightarrow\quad \mathrm{d}y=-e^{-x}\,\mathrm{d}x
\]

For the limits: when \(y=1\), \(e^{-x}=1\), so \(x=0\). The lower limit is improper, so write it as \(y=s\) and let \(s\to 0^{+}\). Then \(x=-\ln s\), and as \(s \to 0^{+}\), \(-\ln s \to \infty\). So the lower limit \(y \to 0^{+}\) corresponds to \(x \to \infty\).

Hence
\begin{align*}
\int_{0}^{1} (\ln y)^2\,\mathrm{d}y &= \int_{\infty}^{0} x^2\left(-e^{-x}\right)\mathrm{d}x\\
&= \int_{0}^{\infty} x^2 e^{-x}\,\mathrm{d}x
\end{align*}

as required, where reversing the limits removes the minus sign.

    \item The integral is defined as
\[
\int_{0}^{\infty} x^2 e^{-x}\,\mathrm{d}x=\lim_{t\to\infty}\int_{0}^{t} x^2 e^{-x}\,\mathrm{d}x
\]

First find the indefinite integral using integration by parts:
\[
\int u\frac{\mathrm{d}v}{\mathrm{d}x}\,\mathrm{d}x = uv - \int v\frac{\mathrm{d}u}{\mathrm{d}x}\,\mathrm{d}x
\]

Take \(u=x^2\), \(\frac{\mathrm{d}v}{\mathrm{d}x}=e^{-x}\), so \(\frac{\mathrm{d}u}{\mathrm{d}x}=2x\), \(v=-e^{-x}\):
\[
\int x^2e^{-x}\,\mathrm{d}x=-x^2e^{-x}+\int 2xe^{-x}\,\mathrm{d}x
\]

For the remaining integral take \(u=2x\), \(\frac{\mathrm{d}v}{\mathrm{d}x}=e^{-x}\), so \(\frac{\mathrm{d}u}{\mathrm{d}x}=2\), \(v=-e^{-x}\):
\begin{align*}
\int 2xe^{-x}\,\mathrm{d}x&=-2xe^{-x}+\int 2e^{-x}\,\mathrm{d}x\\
&=-2xe^{-x}-2e^{-x}+c
\end{align*}

Therefore
\[
\int x^2e^{-x}\,\mathrm{d}x=-e^{-x}\left(x^2+2x+2\right)+c
\]

So
\begin{align*}
\int_{0}^{t} x^2 e^{-x}\,\mathrm{d}x&=\left[-e^{-x}\left(x^2+2x+2\right)\right]_{0}^{t}\\
&=-e^{-t}\left(t^2+2t+2\right)-\left(-e^{0}\times 2\right)\\
&=2-t^2e^{-t}-2te^{-t}-2e^{-t}
\end{align*}

As \(t\to\infty\), the exponential dominates any power of \(t\), so \(t^2e^{-t}\to 0\) and \(te^{-t}\to 0\); also \(e^{-t}\to 0\). Hence
\[
\int_{0}^{\infty} x^2 e^{-x}\,\mathrm{d}x=\lim_{t\to\infty}\left(2-t^2e^{-t}-2te^{-t}-2e^{-t}\right)=2
\]

    \item By part (b)(i) the two integrals are equal, so by part (b)(ii)
\[
\int_{0}^{1} (\ln y)^2\,\mathrm{d}y=2
\]
  \end{enumerate}
\end{enumerate}
\end{djsSolution}

\end{enumerate}
\end{document}
